\chapter{How to Approach the Exam (2022)}

\section{The Sociology Exam Structure}

\begin{KeyConcept}
\begin{itemize}
\item \textbf{Paper structure}: each paper (Paper 1 and Paper 2) has \textbf{28 questions}, of which you \textbf{attempt 19} --- so \textbf{38 questions in six hours}, with a real \textbf{choice}. Unlike GS (where all questions are compulsory), in sociology \textbf{your choice determines your marks}, so the ability to pick the right question is decisive.
\item \textbf{Paper 1}: \textbf{Unit 4 (thinkers) and Unit 5 (stratification)} together carry about \textbf{170 marks out of 250} --- so make notes on these two. \textbf{Paper 2}: \textbf{Section A} (perspectives) and \textbf{Section B} (institutions) are compulsory/static; \textbf{Section C} is contemporary/current-affairs --- except \textbf{social movements, 'vision of social change', and rural/agrarian transformation programmes} (community development, Green Revolution), which are static.
\item \textbf{Sources} for the GS overlap: \textbf{Ram Ahuja} for social issues/social justice; current affairs for the rest.
\end{itemize}
\end{KeyConcept}

\begin{QuickRevision}
\begin{itemize}
\item Paper structure
\item 28 questions
\item attempt 19
\item 38 questions in six hours
\item choice
\item your choice determines your marks
\item Paper 1
\item Unit 4 (thinkers) and Unit 5 (stratification)
\item 170 marks out of 250
\item Paper 2
\item Section A
\item Section B
\item Section C
\item Sources
\end{itemize}
\end{QuickRevision}

\section{Answer-Writing and Essay Guidance}

\begin{KeyConcept}
\begin{itemize}
\item \textbf{GS1 social issues and social justice} largely revolve around \textbf{secularism, communalism, globalisation, and the women's movement} --- cover these dimensions (do \textbf{not} write thinkers here, unlike in sociology). Example: 'the women's movement in India did not address issues of lower-strata women' (2018) is answered through \textbf{intersectionality}.
\item \textbf{Essay writing} is about \textbf{flow, not dimensions or facts} --- the evaluator is a \textbf{general reader}, not a subject expert; you are judged on whether you convey your message clearly. (If the optional and ethics are ready, the essay is almost automatically ready --- only the \emph{art} of writing it needs practice.)
\item \textbf{Writing tests is essential}: committing mistakes and correcting them is how the understanding gets fixed ('befriending the mistakes'); a person who has only watched videos/handouts but not written tests cannot be said to know the subject. Make your \textbf{own list of thinkers} during revision, rather than memorising one.
\end{itemize}
\end{KeyConcept}

\begin{QuickRevision}
\begin{itemize}
\item Paper 1 \& 2: 28 questions, attempt 19 each (38 in 6 hrs); choice decides marks.
\item Paper 1 Unit 4+5 (thinkers + stratification) = $\sim$170 marks; Paper 2 Section A/B compulsory, Section C contemporary.
\item GS1 social issues: secularism, communalism, globalisation, women's movement (no thinkers; e.g. intersectionality).
\item Essay = flow, not dimensions; write tests (befriend mistakes); make your own thinkers list.
\end{itemize}
\end{QuickRevision}
